\input{preamble}
\newpage

\section{Linear regression}

Let $Y$ be a response vector with $n$ rows.
Let $p$ be the number of model parameters.
Then for design matrix $X$ with $n$ rows and $p$ columns the ordinary least squares
estimate of model parameters $\beta$ is
\begin{equation*}
\beta=(X^TX)^{-1}X^TY
\end{equation*}

Example.
Response variable $Y$ is log of half-life for uranium ($Z=90$)
and thorium ($Z=88$).
Note: $Z$ is atomic number of daughter nucleus.
\begin{equation*}
\begin{matrix}
Y & A & Z & E
\\[1ex]
6.30 & 228 & 90 & 6.80
\\
14.4 & 230 & 90 & 5.99
\\
21.5 & 232 & 90 & 5.41
\\
29.7 & 234 & 90 & 4.86
\\
34.2 & 236 & 90 & 4.57
\\
39.5 & 238 & 90 & 4.27
\\
7.52 & 226 & 88 & 6.45
\\
17.9 & 228 & 88 & 5.52
\\
28.5 & 230 & 88 & 4.77
\\
40.6 & 232 & 88 & 4.08
\end{matrix}
\end{equation*}

Geiger–Nuttall regression model.
\begin{equation*}
\hat Y=\beta_1\frac{Z}{\sqrt E}+\beta_2Z^{2/3}+\beta_3
\end{equation*}

The following code has design matrix $X$
as $p\times n$ hence the ordinary least squares formula is
\begin{equation*}
\beta=(XX^T)^{-1}XY
\end{equation*}

{\color{blue}
\begin{verbatim}
X = zero(3,10)
X[1] = Z / sqrt(E)
X[2] = Z^(2/3)
X[3] = X[3] + 1

beta = dot(inv(dot(X,transpose(X))),X,Y)

Yhat = dot(beta,X)
\end{verbatim}}

\href{https://georgeweigt.github.io/examples/linear-regression-demo.html}{Eigenmath script}

\end{document}
